\section{Subgroup and Intersectional Analyses}

\subsection{Subgroups and Stratification Variables to be Examined}

The study will examine subgroups defined by demographic characteristics, fandom history, and patterns of community participation. Potential stratification variables include age, country of residence, educational attainment, political orientation, religious affiliation, LGBTQ+ identity, disability, neurodivergence, generation engagement, fandom-entry cohort, writer/reader status, and level of Fimfiction or broader fandom engagement.

Subgroups may also be defined according to content preferences and community participation, including generation preferences, romance preferences, reading of Mature-rated or sexually explicit fiction, participation in Fimfiction groups, convention attendance, cross-platform engagement, and writer status.

% TBD: The final set of stratification variables and any prespecified primary subgroup classifications
% will be established before analysis.

\subsection{Planned Subgroup Comparisons and Cross-tabulations}

\sourceimage{1695599}{When everypony is more than just the ponies}

The analysis will use cross-tabulations and other descriptive comparisons to examine how characteristics and behaviours vary across subgroups. Planned areas of examination include demographic characteristics by fandom engagement, disability and neurodivergence by LGBTQ+ identity, demographic characteristics by generation engagement, age and fandom-entry characteristics, writer versus reader status, and demographic characteristics by reading and content preferences.

Cross-tabulations will also be used to examine intersections between multiple demographic and fandom characteristics where sample sizes permit. These comparisons will describe differences within the observed respondent population rather than establish causal relationships.

% TBD: The specific cross-tabulations to be reported, the prioritisation of comparisons, and whether formal
% statistical tests of subgroup differences will be performed will be established before analysis.

\subsection{Group-specific Descriptive Statistics}

For sufficiently large subgroups, descriptive statistics will be calculated using the same general framework as for the overall sample. Categorical variables will generally be reported as frequencies and percentages, while numerical and ordinal variables will receive appropriate measures of central tendency, dispersion, and distribution.

Intersectional groups may be examined where meaningful, including combinations of demographic and fandom characteristics.

% TBD: The specific subgroup-specific statistics and visualisations to be reported will be established before analysis.

\subsection{Co-occurrence of Responses or Categories}

Where respondents can legitimately select multiple responses, the study will examine the co-occurrence of those categories. This will be relevant to select-all-that-apply measures of disability, neurodivergence, LGBTQ+ identity, other demographic characteristics, fandom platforms, other fandoms, and potentially content preferences.

Co-occurrence may be described using counts, percentages, cross-tabulations, and visualisations showing the extent to which categories overlap.

% TBD: The specific co-occurrence analyses and the maximum number of dimensions examined simultaneously
% will be established before analysis.

\subsection{Association Statistics, Where Applicable}

Association statistics may be used where they provide an appropriate summary of the relationship between variables. The specific statistic will depend on the measurement characteristics and distribution of the variables being examined.

Any association statistics will be interpreted consistently with the descriptive and exploratory purpose of the study and will not be treated as evidence of causation.

% TBD: Whether formal association statistics will be included, which statistics will be used, and whether
% statistical significance testing or confidence intervals will accompany them will be established before
% those analyses are conducted.

\subsection{Rules for Handling Small Subgroup Sizes}

Small subgroups will be handled cautiously because estimates based on very few respondents may be unstable, potentially identifiable, or misleading. Small cell counts will not be interpreted as reliable estimates of the broader Fimfiction community. Where necessary, categories may be combined into substantively meaningful broader categories rather than reporting highly granular groups with very small numbers.

Small numbers alone will not constitute grounds for excluding participants.

% TBD: The minimum subgroup size required for reporting, suppression thresholds, rules for combining categories,
% and whether exact counts or percentages will be withheld below specified thresholds will be established
% before subgroup analyses are conducted.

\subsection{Statistical Thresholds and Multiple Comparisons}

The primary subgroup and intersectional analyses are descriptive and will not depend on inferential statistical hypothesis testing. Descriptive differences between groups will be interpreted on the basis of their magnitude, distribution, uncertainty where applicable, and substantive relevance.

No statistical significance threshold or correction for multiple comparisons is planned for the primary descriptive analyses.

% TBD: If formal inferential testing is included as a supplementary or exploratory analysis, the relevant
% significance threshold and multiple-comparison procedure will be specified before those analyses are conducted.
